\section{Experiments}
\label{sec:experiments}

We organize the evidence into four questions: the accuracy and
finite-horizon reliability of local prediction, the benefit of
cross-regime routing and fusion, the behavior of internal routing
and physical-field diagnostics, and the computational cost.
The three benchmarks provide complementary rather than identical
tests: Circular H/P supplies controlled local comparisons, Square
supplies the principal history-conditioned fusion study, and Pinball
provides an archived POD--Galerkin comparison and additional overlap
results. Chart-wise parameter ranges, reduced dimensions, and trajectory
splits are summarized in Table~\ref{tab:chart_splits}.
Evaluation horizons are specified separately for each comparison.

Relative weighted velocity and gauge-corrected pressure errors
$E_u,E_p$ are percentages, with $E_{\rm joint}=E_u+E_p$.
The new Circular controls and Square reevaluation use
POD-reconstructed reference fields; their errors exclude the
projection remainder relative to raw CFD.
Unless otherwise stated for these reevaluations, errors are averaged
over prediction steps, then windows within each Reynolds number,
then equally over Reynolds numbers.
We distinguish training failures, validation rejection, and
non-finite test predictions. Error statistics over accepted runs
alone are not complete multi-seed robustness estimates.
Appendix~\ref{app:implementation} specifies metrics and data splits.



\subsection{Local prediction accuracy and finite-horizon reliability}
\label{sec:exp_local}

\paragraph{Comparison with POD--Galerkin.}
A local predictor should remain accurate after reference-state
injection stops. Table~\ref{tab:pinball_autonomous} retains the
archived Pinball comparison with classical POD--Galerkin over
$K=32$--$56$ steps. The Periodic joint errors are $0.8415\%$
at $K=32$ and $1.672\%$ in the $K=56$ periodic-core evaluation,
compared with $43.20\%$ and $41.91\%$ for POD--Galerkin.
The periodic-core cadence is $\Delta t=1$, giving 56 physical
time units of prediction.

These values demonstrate lower reported errors over the evaluated
windows, not arbitrary-time stability. Moreover, the Periodic
result chain identifies Deep-FNN-H3: its learned velocity derivative
uses Galerkin information as a feature rather than adding a
sparse-MoE correction to the Galerkin output.
Standalone Steady/Hopf result identities remain unresolved.
We therefore do not interpret this archived table as an isolated
test of adding sparse MoE to POD--Galerkin.
A matched bare-Galerkin/additive-MoE comparison with the same chart,
initialization, time stepping, pressure treatment, and test windows
is not established by the current evidence.

\begin{table}[H]
\centering
\caption{
Archived autonomous local-predictor comparison on fluidic pinball. The Periodic values are associated with Deep-FNN-H3, not an additive sparse-MoE correction; standalone Steady/Hopf identities remain unresolved. Errors are physical-field percentages; lower is better. Periodic-core denotes the longer-horizon periodic evaluation; complete protocol details are given in Appendix~\ref{app:pinball}. A dash indicates that no baseline error is reported.
}
\label{tab:pinball_autonomous}
\small
\setlength{\tabcolsep}{5pt}
\begin{tabular}{@{}llrrr@{}}
\toprule
Evaluation setting & Model & $E_u$ & $E_p$ & $E_{\mathrm{joint}}$ \\
\midrule
Steady ($K=56$)
& POD--Galerkin
& -- & -- & -- \\
& Reported local predictor
& 0.2247 & 0.7593 & 0.9840 \\
\addlinespace

Hopf ($K=56$)
& POD--Galerkin
& 0.4378 & 48.50 & 48.94 \\
& Reported local predictor
& 0.0995 & 1.986 & 2.085 \\
\addlinespace

Periodic ($K=32$)
& POD--Galerkin
& 1.641 & 41.56 & 43.20 \\
& Reported local predictor
& 0.2897 & 0.5518 & 0.8415 \\
\addlinespace

Periodic-core ($K=56$)
& POD--Galerkin
& 2.555 & 39.36 & 41.91 \\
& Reported local predictor
& 0.6222 & 1.050 & 1.672 \\
\bottomrule
\end{tabular}
\end{table}

\paragraph{Controlled local comparisons.}
Circular H/P is reevaluated on common held-out windows with seeds
1248, 1600, and 2026. Dense matches the active neural parameter
budget; Structured retains explicit response branches but removes
input-dependent expert routing. These controls answer a different
question from the bare POD--Galerkin comparison.
Table~\ref{tab:circular_full_reorg} reports Periodic $K=24$ results;
$K=48$, Hopf results, the OpInf adaptation, and acceptance counts
are provided in Appendix~\ref{app:circular}.

\begin{table}[H]
\centering
\caption{Circular Periodic controls: full-window mean physical-field
errors (\%) at $K=24$, against POD-reconstructed truth. Entries are mean
$\pm$ sample standard deviation over seeds 1248, 1600, and 2026; all
listed neural configurations have 3/3 accepted seeds. Repaired denotes
the changed quadratic-branch initialization.}
\label{tab:circular_full_reorg}
\small
\setlength{\tabcolsep}{4pt}
\begin{tabular}{@{}lccc@{}}
\toprule
Model & $E_u$ & $E_p$ & $E_u+E_p$\\
\midrule
Proposed, original & $0.5731\pm0.3638$ & $2.3060\pm1.3681$ & $2.8791\pm1.7269$\\
Proposed, repaired & $0.5814\pm0.3949$ & $2.4854\pm1.4427$ & $3.0667\pm1.8373$\\
Dense, active-matched & $\mathbf{0.2446\pm0.0323}$ & $\mathbf{1.0805\pm0.1946}$ & $\mathbf{1.3251\pm0.2268}$\\
Structured, original & $0.5178\pm0.2667$ & $2.3489\pm1.1688$ & $2.8666\pm1.4342$\\
Structured, repaired & $0.5429\pm0.2783$ & $2.3168\pm1.0315$ & $2.8598\pm1.3081$\\
\bottomrule
\end{tabular}
\end{table}

Dense is more accurate than the proposed configurations at both
horizons, while Structured is close to the original proposed model.
Thus these results do not identify sparse routing as a Periodic
accuracy advantage under the matched active-parameter budget.
The original quadratic factors remained zero after joint zero
initialization. Repairing the initialization activates the branch
but does not consistently improve accuracy.
Only two of three repaired FP32 Hopf runs per architecture are
accepted; their conditional errors do not establish complete
three-seed robustness.
The task-specific OpInf fit is competitive on Hopf
($0.03870\%$ joint error at $K=24$), but poor on Periodic
($114.0941\%$); its boundary-selected regularization limits a
general conclusion about that method family.

\subsection{Cross-regime routing and physical-space fusion}
\label{sec:exp_fusion}

We hold the Square local predictors and E2 router fixed and compare
assembly rules on the prescribed S--H and H--P support at $K=24$.
Each interface contains 16 held-out windows.
E2 Top-1 selects within the supplied pair; T2-C learns a correction
to its pair preference using initial-history descriptors.
Neither the candidate set nor an automatic applicability mask is
learned in this experiment.

Table~\ref{tab:routing_fusion} uses a common full-window metric.
Train-fixed $\alpha$ and train-fixed logit offset are fitted only
on training windows, and the fixed candidate is selected on
validation error. These controls test whether the improvement
can be explained by a constant mixture or calibration offset.

\begin{table}[H]
\centering
\small
\caption{Square full-window joint error (\%) at $K=24$ against the
cached POD reference. Entries use the frozen-checkpoint reevaluation,
not a three-seed average. Scalar controls use training data;
the fixed candidate uses validation data.}
\label{tab:routing_fusion}
\begin{tabular}{@{}lrr@{}}
\toprule
Assembly rule & S--H & H--P\\
\midrule
E2 Top-1 & 1.5786 & 1.6424\\
E2-probability blend & 1.8145 & 1.4489\\
Equal blend & 2.3849 & 1.3930\\
Train-fixed $\alpha$ & 2.0591 & 0.9767\\
Train-fixed logit offset & 2.0149 & 0.9821\\
Validation-fixed candidate & 1.5786 & 1.0068\\
T2-C & \textbf{1.0954} & \textbf{0.8123}\\
\bottomrule
\end{tabular}
\end{table}

T2-C improves on all listed controls for these fixed candidate
sets. A separate architecture-matched $\mu$-only control zeros the
standardized history channels while retaining the gate architecture.
Across three paired gate seeds, full-window joint errors are
$1.095351\pm0.000015\%$ versus $1.814699\pm0.000242\%$
in S--H and $0.812343\pm0.000177\%$ versus
$0.976939\pm0.000179\%$ in H--P.
The reductions are $39.64\%$ and $16.85\%$.
These quantify gate-training variability, not uncertainty from
retraining the complete hierarchy.

\begin{figure}[H]
\centering
\input{figures/t2c_window_comparison.tex}
\caption{Full-window Square comparison over three paired gate seeds.
All local candidates and E2 are fixed; bars show mean and sample
standard deviation.}
\label{fig:square_overlap}
\end{figure}

The evidence supports history-conditioned fusion for these windows,
but does not prove that history dependence is necessary for every
flow or gate. The weight is fixed within each prediction window;
history can change it between windows.
Appendix~\ref{app:square} retains component errors, terminal-step
results, scalar objectives, and complete field visualizations.
Additional Pinball overlap results appear in
Appendix~\ref{app:pinball_overlap_results}; their local candidates
are Deep-FNN-H3, so they test the assembly interface rather than
the sparse-MoE architecture.
A posteriori convex oracles are diagnostics, not deployable methods.

\subsection{Internal routing and physical-field diagnostics}
\label{sec:exp_diagnostics}

\paragraph{Expert utilization.}
Post-hoc diagnostics use frozen Circular H/P specialists and are
not used for checkpoint selection.
Periodic distributes decisions across group--pair combinations;
Hopf has a prescribed single group and fixed sparse support, but
nonuniform, parameter-dependent mixture weights.
Its single group is not evidence of collapse of a learned
multi-group router. The detailed frequency, weight-mass, and entropy
statistics in Appendix~\ref{app:internal_routing} count the first
integration stage per output step, not all RK evaluations.
Utilization alone does not establish an accuracy benefit.

\paragraph{Fusion consistency.}
Appendix~\ref{app:fusion_consistency} separates conditional algebraic
identities from measured reconstructed-field diagnostics.
On both Square interfaces, the measured T2-C pressure residual is
lower than either candidate for both tested spatial stencils.
However, Hopf has a smaller H--P momentum residual and a smaller
S--H energy error. Fusion is therefore not a general residual or
energy-error minimizer.

The measurements use offline least-squares spatial derivatives
and coarse-time differences, not the original OpenFOAM
finite-volume residuals. Nonzero reference residuals include POD
truncation and differentiation error.
Identity discrepancies near machine precision verify the algebra,
not conservation or physical accuracy.
The phase screen resolves none of the 16 S--H windows and 15 of
16 H--P windows; even resolved short-window traces do not establish
absence of aliasing or long-time phase stability.
The complete energy, residual, probe, and phase curves are retained
in Appendix~\ref{app:fusion_measured}.

\subsection{Computational cost and timing boundaries}
\label{sec:exp_cost}

Sparse activation is a parameter-use property, not a runtime claim.
Circular H/P activate $44.53\%$ and $17.09\%$ of their trainable
parameters per analyzed prediction decision.
For native Circular Periodic prediction at batch size one and
$K=48$, median times are 2.1848 s (proposed), 0.6579 s (Dense),
and 0.8652 s (Structured).
The proposed implementation is therefore slower than Dense in this
measurement despite sparse activation; the native Galerkin path
also executes on the CPU.

A separate warm physical-input/physical-output measurement includes
history encoding, local advancement, and all-step field decoding.
It takes 2.18747 s on an RTX 4090 with a Xeon host for a Circular
Periodic window spanning 598.70093 physical time units.
A matched historical production FOM interval takes 1996 s.
The resulting ratio of approximately 912.5 is not a full-hierarchy
or hardware-independent speedup: hardware and I/O boundaries differ,
the FOM timing has one realization, and ROM loading, output writes,
E2, descriptors, and T2-C are excluded.
Appendix~\ref{app:parameter_counts} lists warm-ups, repetitions,
synchronization, memory, hardware, and all included/excluded work.
A fully paired Top-2 end-to-end timing remains unmeasured.
